\specialchapter{Index of notation}

The reference numbers indicate respectively the chapter, paragraph and number (or, occasionally, exercise).

\([x,y]\) \((x,y\) elements of a Lie algebra): I.1.2.

\[
\mathfrak {g} ^ {0}
\]

\[
\mathfrak {g l} (\mathrm{E}), \mathfrak {g l} (n, \mathrm{K})
\]

\([\mathfrak{a},\mathfrak{b}],[z,\mathfrak{a}],[\mathfrak{a},z]\) (a, b submodules, \(z\) an element of a Lie algebra): I.1.4.

af(M) (M a K-module): I.1.8.

\(U_{+}, U_{0}\) (U the enveloping algebra of a Lie algebra): I.2.1.

\(x_{\mathbb{M}}\) ( \(x\) an element of a Lie algebra \(\mathfrak{g}\) , M a \(\mathfrak{g}\) -module): I.3.1.

\[
\mathrm{T} _ {n}, \mathrm{U} _ {n}, \mathrm{G} ^ {n}: \text { I.2.6. }
\]

\(C(\rho)\) ( \(\rho\) a representation of a Lie algebra): I.7.1.

\[
\mathcal {C} ^ {\infty} \mathfrak {g}, \mathcal {D} ^ {\infty} \mathfrak {g}: \mathrm{I}. 1, \text { Exercise 14. }
\]

\(x^{[p]}\) : I.1, Exercise 20.

\(\mathbf{GL}(n,\mathbf{K})\) (formal group): I.1, Exercise 25.

\(\mathfrak{o}(\Phi)\) : I.1, Exercise 26.

\[
\mathrm{C} ^ {p} (\mathfrak {g}, \mathrm{M}), \mathrm{C} ^ {*} (\mathfrak {g}, \mathrm{M}), i (y), \theta (x), d, \mathrm{Z} ^ {p} (\mathfrak {g}, \mathrm{M}), \mathrm{B} ^ {p} (\mathfrak {g}, \mathrm{M}), \mathrm{H} ^ {p} (\mathfrak {g}, \mathrm{M}), \mathrm{H} ^ {*} (\mathfrak {g}, \mathrm{M})
\]

\(\mathfrak{sp}(2n,\mathbf{K})\) : I.6, Exercise 25.

K: II. Conventions.

g, U = Ug, σ: g → Ug: II.1.

E, E\(^{+}\): II.1.1. P(E), π, γ, c\(^{+}\): II.1.6. S(g), ε, γ: II.1.5. f\(_{E}\): U(P(E)) → E: II.1.6. M(X), l(w), Lib(X) = Lib\(_{K}\)(X): II.2.1. L(X) = L\(_{K}\)(X): II.2.2. φ: X → L(X): II.2.2. (a, r): II.2.3. L(u): II.2.5. Lib\(^{b}\)(X), L\(^{b}\)(X), L\(^{n}\)(X): II.2.6. P\(_{n}\): II.2.7. C\(^{n}\)(g): II.2.7. H, d\(_{y}\): II.2.10. w = Ψ(w): II.2.11. A(X) = A\(_{K}\)(X), A\(^{+}\)(X), Mo(X): II.3. π: II.3.2. (G\(_{\alpha}\)), (G\(_{\alpha}\) \(^{+}\)): II.4.1. v: II.4.2. gr(G), gr\(_{\alpha}\)(G): II.4.3. F(X), A(X), A\(^{n}\)(X): II.5. Â(X), ω: II.5.1. ε(a): II.5.2. l(x), exp(x), log(y): II.6.1. e(X), l(X): II.6.1. L(X): II.6.2. a ⊕ b: II.6.2. H, H\(_{n}\), H\(_{ra}\): II.6.4. H, Ω: II.7.2. A, exp\(_{A}\), log\(_{A}\), P(A\(^{I}\), A): II.7.3. v, θ = \(\frac{1}{p - 1}\): II.8. S(n): II.8.1. h(x, y): II.8.3. G\(_{R}\): II.8.4. μ(n): II.App. e, e\(_{G}\), γ(g), δ(g), Int(g), f: III.Conventions. GL(E), GL(n, K): III.1.1, III.3.10. G\(^{v}\): III.1.2. τ(g), ρ(x): III.1.5. (G, g, θ, m): III.1.10. T(m): III.2.1. T(G), T(φ): III.2.2. t * t': III.3.1, III.3.18.

U(G), U \(^{+}\) (G), U \(_{a}\) (G), U \(^{+}\) (G): III.3.1, III.3.18. T \(^{(y)}\) (G), T \(^{(\omega)}\) (G), T \(^{(\infty)}\) (G): III.3.1, III.3.18. t * f: III.3.4, III.3.18. D \(_{i}\) : III.3.5, III.3.18. L \(_{t}\) , R \(_{t}\) : III.3.6, III.3.18. L(G): III.3.7, III.3.18. L(φ): III.3.8, III.3.18. \textless{} t,f\textgreater{} : III.3.9, III.3.18. SL(E): III.3.10. Ad, Ad(g) : III.3.12. {[}α{]} \(^{2}\) : III.3.14. mod(ω) \(_{u}\) , mod φ: III.3.16. f \(^{-1}\) . df: III.3.17, III.3.18. H: III.4.2. g \(^{t}\) , ω(x) : III.4.3. x.y, x{[}n{]}: III.5. c \(_{\alpha\beta y}\) , B(x,y) : III.5.1. e \(_{\alpha}\) : III.5.2. ψ \(_{j}\) , ψ \(_{p,m}\) , \(\binom{l}{i}\) : III.5.3. E(x), L(x) : III.5.4. P \(_{M,x,y}\) : III.6.2. Ad(a) = Int(a): III.6.2. exp, exp \(_{G}\) : III.6.4. Ad(G) = Int(L(G)): III.6.4. L(ρ): III.6.5. G: III.6.10. A, m, p: III.7. G(a): III.7.4. h \(_{n}\) : III.7.5. G \(_{f}\) , log \(_{G}\) , log: III.7.6. D \(^{i}\) G, C \(^{i}\) G: III.9.1. Z \(_{G}\) (A), Z \(_{G}\) (a), λ \(_{g}\) (A), λ \(_{g}\) (a): III.9.3. N \(_{G}\) (A), N \(_{G}\) (a), n \(_{p}\) (a): III.9.4. R, N, r, n: III.9.7. π \(_{1}\) ⊗⋯⊗ π \(_{n}\) , T(π), S(π), ∧(π), T \(^{n}\) (π), S \(^{n}\) (π), ∧ \(^{n}\) (π): III.App.

\specialchapter{Index of terminology}

The reference numbers indicate respectively the chapter, paragraph and number (or, occasionally, exercise)

Adjoint group of a real or complex Lie group: III.6.4. linear mapping of an element of a Lie algebra: I.1.2. representation of a Lie algebra: I.3.1. representation of a Lie group: III.3.12.

Ado's Theorem: I.7.3.

Algebra (not necessarily associative): I.1.1. derived from an algebra by extending the scalars: I.1.1. enveloping, of a Lie algebra: I.2.1. Lie: I.1.2. opposite: I.1.1. product: I.1.1. quotient: I.1.1. restricted enveloping, of a Lie p-algebra: I.2, Exercise 6. symmetric, of a module: I.2.5.

Almost simple Lie group: III.9.8.

Alternants of degree n: II.2.6.

Antiautomorphism, principal, of the enveloping algebra of a Lie algebra: I.2.4.

Associated (bilinear form) with a g-module (a representation): I.3.6. (law of infinitesimal operation) with a law of operation: III.3.7.

Associative algebra, free: II.3.

Automorphism, special, of a Lie algebra: I.6.8.

Basic commutators: II.5.4. family of a Lie algebra: II.2.3.

Bieberbach's Theorem: III.4, Exercise 13.

Bigebra: II.1.2.

Biinvariant section: III.3.13.

Binomial polynomial: II.5, Exercise 4.

Bracket: I.1.2.

Canonical differential form, left: III.3.13, III.3.18. Cartan, criterion: I.5.4. Casimir element: I.3.7. Central extension: I.1.7. filtration of a group: II.4.4. Centralizer: I.1.6, III.9.3. Centre of a Lie algebra: I.1.6. Characteristic ideal: I.1.4. Class, nilpotency: II.2.7. of simple representations: I.3.1. Coboundaries, cochains, cocycles with values in a g-module: I.3, Exercise 12. Cogebra: II.1.1. Commutative Lie algebra: I.1.3. Compatible group and manifold structures: III.1.1. Completely invariant bilinear form: I.3.6. reducible representation: I.3.1. Complex Lie group: III.1.1, III.8.1. Complexification of a real Lie group: III.6.10. Component, isotypical, of a g-module: I.3.1. simple, of a semi-simple Lie algebra: I.6.2. Conjugate Lie group of a complex Lie group: III.1.1. Constant term of an element of the enveloping algebra of a Lie algebra: I.2.1. term of an element of the Magnus algebra: II.5.2. Constants of structure of an algebra with respect to a basis: I.1.1. Containing a simple representation n times (representation): I.3.1. Contragredient of an analytic representation: III.3.11. Convolution of a point distribution and a function: III.3.4. product: III.3.1, III.3.18. Counit of a cogebra: II.1.1. C'-connected subset of a Lie group: III.6.2. Criterion, Cartan: I.5.4.

Derivation of an algebra: I.1.1. inner, of a Lie algebra: I.1.2. left invariant, relative to a formal group: I.1, Exercise 24. Derivative, partial, in the algebra of a free group: II.5, Exercise 2. Derived (algebra) from an algebra by extending the scalars: I.1.1. (representation) from a representation by extending the scalars: I.3.8. ideal: I.1.5. series: I.1.5.

Differential, left, of a mapping into a Lie group: III.3.17, III.3.18. Dimension of a representation: I.3.1. Dual representation: I.3.3.

Element, Casimir: I.3.7. nilpotent: I.6.3. of degree n in the graded algebra associated with the enveloping algebra of a Lie algebra: I.2.6. of filtration \(\leqslant n\) in the enveloping algebra of a Lie algebra: I.2.6. semi-simple: I.6.3. Elements, permutable, in a Lie algebra: I.1.3. Elimination theorem: II.2.9. Engel's Theorem: I.4.2. Enveloping algebra of a Lie algebra: I.2.1. bigebra: II.1.4. Equivalent extensions of a Lie algebra: I.1.7. Exhaustive filtration: II.4.1. Exponential: II.6.1, III.4.3, III.6.4. Extension, central: I.1.7. inessential: I.1.7. of a g-module, of a representation of g: I.7.2. of a Lie algebra by a Lie algebra: I.1.7. trivial: I.1.7. Extensions, equivalent: I.1.7.

Faithful representations: I.3.1. Field of point distributions: III.3.5, III.3.18. Filtered bigebra: II.1.3. Filtration, real, on a group: II.4.1. First species, canonical chart of the: III.4.3. species, system of canonical coordinates of the: III.4.3. Foliation, left, associated with a Lie subalgebra: III.4.1. Form, bilinear, associated with a g-module (a representation of g): I.3.6. bilinear, completely invariant: I.3.6. bilinear, invariant: I.3.6. Killing: I.3.6. Formulae, Jacobson: I.1, Exercise 19. Free Lie algebra: II.2.2. Lie p-algebra: II.3, Exercise 4. magma: II.2.1. Function, order, associated with a filtration: II.4.2.

Graded group associated with a filtered group: II.4.3. Lie algebra, associated: II.4.4. Graduation, total, of L(I): II.2.6. Group commutative formal: I.1, Exercise 24. formal: I.1, Exercise 24. formal linear: I.1, Exercise 25. formal orthogonal: I.1, Exercise 26. formal symplectic: I.1, Exercise 26. germ, Lie: III.1.10. germ, Lie, defined by a Lie algebra: III.4.2.

Hall basis: II.2.11. formula: II.5, Exercise 9. set: II.2.10. Hausdorff function: II.7.2, II.8.3. group: II.6.2. inversion formula: II.6, Exercise 4. series: II.6.4. Holomorph of a Lie algebra: I.1, Exercise 16. Homomorphism, algebra: I.1.1. canonical, of the enveloping algebra of a Lie subalgebra of g into the enveloping algebra of g: I.2.3. canonical, of the symmetric algebra of g onto the graded algebra associated with the enveloping algebra of g: I.2.6. formal: I.1, Exercise 24.

Ideal, characteristic: I.1.4. derived: I.1.5. left (right, two-sided) of an algebra: I.1.1. of a Lie algebra: I.1.4. Identity, Jacobi: I.1.2. Induced Lie group structure: III.4.5. Inessential extension: I.1.7. Infinitesimal automorphism: III.10.1. Inner derivation: I.1.2. Integral filtration of a group: II.4.1. subgroup of a Lie group: III.6.2. Invariant bilinear form: I.3.6. left, point distributions, field of: III.3.6. of a g-module (a representation): I.3.5. section: III.3.13.

Inverse image Lie group structure: III.1.9. Irreducible representation: I.3.1. Isomorphic representations: I.3.1. Isomorphism, canonical, of the symmetric algebra of g onto the underlying vector space of its enveloping algebra: I.2.7. Isotypical component of a g-module: I.3.1.

Jacobi identity: I.1.2. Jordan's Theorem: III.4, Exercise 11.

Kernel of an extension: I.1.7. Killing form: I.3.6.

Largest nilpotency ideal of a g-module (a representation): I.4.3. Largest nilpotent ideal of a Lie algebra: I.4.4. Law chunk of operation: III.1.11. formal group: I.1, Exercise 24. of infinitesimal operation: III.3.7, III.3.18. Length of an element of a free magma: II.2.1. Levi subalgebra: I.6.8. Levi-Malcev Theorem: I.6.8. Lie algebra: I.1.2. algebra, commutative: I.1.3. algebra, characteristically nilpotent: I.4, Exercise 19. algebra, nilpotent: I.4.1. algebra of a formal group: I.1, Exercise 24. algebra of a Lie group: III.3.7. algebra of a Lie group germ: III.3.18. algebra of endomorphisms of a module: I.1.2. algebra, reductive: I.6.4. algebra, semi-simple: I.6.1. algebra, simple: I.6.2. algebra, solvable: I.5.1. group: III.1.1. Lie's Theorem: I.5.3. Locally isomorphic group germ: III.1.10. Logarithm: II.6.1, III.7.6. Lower central series: I.1.5.

Magnus algebra: II.5.1. group: II.5.2. Mapping, canonical, of a Lie algebra into its enveloping algebra: I.2.1. Maurer-Cartan formulae: III.3.14, III.3.18.

Möbius function: II, Appendix. inversion formula: II, Appendix. Morphism, Lie group: III.1.2. Lie group germ: III.1.10. Multigraduation of L(I): II.2.6.

Nilpotent Lie algebra: I.4.1, II.2.7. radical of a Lie algebra: I.5.3. Normable algebra: II.7. Normalizer: I.1.4, III.9.4. Normed Lie algebra: II.7, II.8.2.

Opposite algebra: I.1.1. Order of an element in a filtered group: II.4.2.

p-adic Lie group: III.1.1, III.8.1. p-algebra, Lie: I.1, Exercise 20. p-unipotent Lie: I.4, Exercise 23. p-core of a commutative Lie p-algebra: I.1, Exercise 23. p-derivation: I.2, Exercise 7. P-envelope of a nilpotent group: II.4, Exercise 15. p-homomorphism: I.1, Exercise 20. p-ideal: I.1, Exercise 22. P-integer: II.4, Exercise 14. p-mapping: I.1, Exercise 20. p-polynomial: I.7, Exercise 5. P-saturation of a subgroup of a nilpotent group: II.4, Exercise 14. P-torsion group, P-torsion-free group: II.4, Exercise 14. Permutable elements: I.1.3, III.9.3. Poincaré-Birkhoff-Witt Theorem: I.2.7. Polynomial, Lie: II.2.4. mapping: II.2.4. Presentation of a Lie algebra: II.2.3. Primitive element of a bigebra: II.1.2. Product algebra: I.1.1. of Lie groups: III.1.4. semi-direct, of Lie algebras: I.1.8. tensor, of representations: I.3.2, III, Appendix. Pure g-module of species N: I.3.1. representation: I.3.1.

Quasi-subgroup, Lie: III.1.3.

Quotient algebra: I.1.1. of a Lie group: III.1.6. representation: I.3.1.

Radical, nilpotent, of a Lie algebra: I.5.3. of a Lie algebra: I.5.2. of a Lie group: III.9.7.

Real Lie group: III.1.1, III.8.1, III.8.2.

Reductive Lie algebra: I.6.4. (Lie subalgebra) in a Lie algebra: I.6.6.

Relators: II.2.3.

Replica of an endomorphism: I.5, Exercise 14.

Representation, adjoint: I.3.1. analytic linear, of a Lie group: III.1.2. completely reducible: I.3.1. containing a simple representation n times: I.3.1. dual: I.3.3. faithful: I.3.1. irreducible: I.3.1. obtained by extending the ring of scalars: I.3.8. of a Lie algebra: I.3.1. pure, of species σ: I.3.1. quotient: I.3.1. semi-simple: I.3.1. simple: I.3.1.

Representations, isomorphic: I.3.1. similar: I.3.1.

Restriction of scalars, Lie group derived from a Lie group by: III.1.1.

Roots of a solvable Lie algebra: III.9, Exercise 17.

Second species, canonical chart of the: III.4.3. species, system of canonical coordinates of the: III.4.3.

Section of a vector bundle: III.1.8.

Semi-direct product of Lie algebras: I.1.8. product of Lie groups: III.1.4.

Semi-simple Lie algebra: I.6.1. Lie group: III.9.8. representation: I.3.1.

Separated filtration: II.4.1.

Series, composition, joining two subalgebras: I.1, Exercise 14. derived: I.1.5. Lie formal power: II.6.3. lower central, of a Lie algebra: I.1.5, II.2.7. upper central, of a Lie algebra: I.1.6.

Similar representations: I.3.1. Simple component of a semi-simple Lie algebra: I.6.2. Lie algebra: I.6.2. representation: I.3.1. Solvable Lie algebra: I.5.1. Space, cohomology, with values in a g-module: I.3, Exercise 12. Lie homogeneous: III.1.6. Special automorphism: I.6.8. Species of an isotypical component, of a pure g-module: I.3.1. Standard group: III.7.3. Subalgebra: I.1.1. Lie: I.6.8. reductive in a Lie algebra: I.6.6. subinvariant: I.1, Exercise 14. Subgroup germ, Lie: III.1.10. Lie: III.1.3. Subrepresentation: I.3.1. Sum, direct, of representations: I.3.1. Symmetric algebra of a module: I.2.5.

t-th power mapping: III.4.3. Tangent law of composition: III.2.1. Lie subgroup: III.4.5. Theorem, Ado's: I.7.3. Engel's: I.4.2. Levi-Malcev: I.6.8. Lie's: I.5.3. Poincaré-Birkhoff-Witt: I.2.7. Weyl's: I.6.2. Zassenhaus's: I.7.2. Transporter: I.1, Exercise 1. Triangular group, upper strict: II.4.6. Trivial extension: I.1.7. g-module: I.3.1. vector G-bundle: III.1.8. Trivialization, right (resp. left), of T(G): III.2.1, III.2.2. Two-sided ideal, I.1.1. Type (N), real Lie group of: III.9, Exercise 29.

\(u\) -primitive elements of a cogebra: II.1.1.\\
Unipotent endomorphism: III.9.5.\\
Universal covering of a connected Lie group: III.1.9.

Upper central series of a group: II.4, Exercise 18. central series of a Lie algebra: I.1.6.

Weyl's Theorem: I.6.2.

Zassenhaus's Theorem: I.7.2.

\summarychapter{Summary of certain properties of finite-dimensional Lie algebras over a field of characteristic 0}{Summary}{of certain properties of finite-dimensional Lie algebras over a field of characteristic 0.}


Let g be a Lie algebra, r its radical, n its largest nilpotent ideal, s its nilpotent radical and t the orthogonal of g relative to the Killing form. Then r, n, s, t are characteristic ideals and \(r \supset t \supset n \supset s\) .

\begin{enumerate}
\def\labelenumi{(\Roman{enumi})}
\tightlist
\item
  Any one of the following properties characterizes semi-simple Lie algebras:
\end{enumerate}

\begin{enumerate}
\def\labelenumi{(\arabic{enumi})}
\tightlist
\item
  \(r = \{0\}\) ; (2) \(n = \{0\}\) ; (3) \(t = \{0\}\) ; (4) every commutative ideal of g is zero; (5) the algebra g is isomorphic to a product of simple Lie algebras; (6) every finite-dimensional representation of g is semi-simple.
\end{enumerate}

\begin{enumerate}
\def\labelenumi{(\Roman{enumi})}
\setcounter{enumi}{1}
\tightlist
\item
  Any one of the following properties characterizes reductive Lie algebras:
\end{enumerate}

\begin{enumerate}
\def\labelenumi{(\arabic{enumi})}
\tightlist
\item
  \(\mathfrak{s} = \{0\}\) ; (2) \(\mathfrak{r}\) is the centre of \(\mathfrak{g}\) ; (3) \(\mathcal{D}\mathfrak{g}\) is semi-simple; (4) \(\mathfrak{g}\) is the product of a semi-simple algebra and a commutative algebra; (5) the adjoint representation of \(\mathfrak{g}\) is semi-simple; (6) \(\mathfrak{g}\) has a finite-dimensional representation such that the associated bilinear form is non-degenerate; (7) \(\mathfrak{g}\) has a finite-dimensional semi-simple faithful representation.
\end{enumerate}

\begin{enumerate}
\def\labelenumi{(\Roman{enumi})}
\setcounter{enumi}{2}
\tightlist
\item
  Any one of the following properties characterizes solvable Lie algebras:
\end{enumerate}

\begin{enumerate}
\def\labelenumi{(\arabic{enumi})}
\tightlist
\item
  \(\mathcal{D}^{p}\mathfrak{g} = \{0\}\) for sufficiently large p; (2) there exists a decreasing sequence \(g = g_{0} \supset g_{1} \supset \cdots \supset g_{n} = \{0\}\) of ideals of g such that the algebras \(g_{i}/g_{i+1}\) are commutative; (3) there exists a decreasing sequence
\end{enumerate}

\[
\mathfrak {g} = \mathfrak {g} _ {0} ^ {\prime} \supset \mathfrak {g} _ {1} ^ {\prime} \supset \dots \supset \mathfrak {g} _ {n ^ {\prime}} ^ {\prime} = \{0 \}
\]

of subalgebras of \(\mathfrak{g}\) such that \(\mathfrak{g}_{i+1}^{\prime}\) is an ideal of \(\mathfrak{g}_i^{\prime}\) and \(\mathfrak{g}_i^{\prime}/\mathfrak{g}_{i+1}^{\prime}\) is commutative; (4) there exists a decreasing sequence \(\mathfrak{g} = \mathfrak{g}_0'' \supset \mathfrak{g}_1'' \supset \cdots \supset \mathfrak{g}_{n''}'' = \{0\}\) of subalgebras of \(\mathfrak{g}\) such that \(\mathfrak{g}_{i+1}''\) is an ideal of codimension 1 in \(\mathfrak{g}_i''\) ; (5) \(\mathfrak{g} \supset \mathcal{D}\mathfrak{g}\) ; (6) \(\mathcal{D}\mathfrak{g}\) is nilpotent.

\begin{enumerate}
\def\labelenumi{(\Roman{enumi})}
\setcounter{enumi}{3}
\tightlist
\item
  Any one of the following properties characterizes nilpotent Lie algebras:
\end{enumerate}

\begin{enumerate}
\def\labelenumi{(\arabic{enumi})}
\tightlist
\item
  \(\mathcal{C}_{g}^{p} = \{0\}\) for sufficiently large p; (2) \(\mathcal{C}_{p}\mathfrak{g} = \mathfrak{g}\) for sufficiently large p; (3) there exists a decreasing sequence \(g = g_{0} \supset g_{1} \supset \cdots \supset g_{p} = \{0\}\) of ideals of g such that \([g, g_{i}] \subset g_{i+1}\) ; (4) there exists a decreasing sequence \(g = g_{0}' \supset g_{1}' \supset \cdots \supset g_{p'}' = \{0\}\) of ideals of g such that \([g, g_{i}'] \subset g_{i+1}'\) and the \(g_{i}'/g_{i+1}'\) are 1-dimensional; (5) there exists an integer i such that
\end{enumerate}

\[
\left(\operatorname{ad} x _ {1}\right) \circ \left(\operatorname{ad} x _ {2}\right) \circ \dots \circ \left(\operatorname{ad} x _ {i}\right) = 0
\]

for all \(x_{1},\ldots ,x_{i}\) in \(\mathfrak{g}\) ; (6) for all \(x\in \mathfrak{g}\) , ad \(x\) is nilpotent.

\begin{enumerate}
\def\labelenumi{(\Alph{enumi})}
\setcounter{enumi}{21}
\tightlist
\item
  g commutative \(\Rightarrow\) g nilpotent \(\Rightarrow\) the Killing form of g is zero \(\Rightarrow\) g solvable.
\end{enumerate}

g commutative \(\Rightarrow\) g reductive.

g semi-simple \(\Rightarrow\) g reductive.

\textbf{(VI) Characterizations of r:}

\begin{enumerate}
\def\labelenumi{(\arabic{enumi})}
\tightlist
\item
  r is the largest solvable ideal of g; (2) r is the smallest ideal such that g/r is semi-simple; (3) r is the only solvable ideal such that g/r is semi-simple; (4) r is the orthogonal of \(\mathcal{D}\mathfrak{g}\) relative to the Killing form.
\end{enumerate}

\textbf{(VII) Characterizations of n:}

\begin{enumerate}
\def\labelenumi{(\arabic{enumi})}
\tightlist
\item
  n is the largest nilpotent ideal of g; (2) n is the largest nilpotent ideal of r; (3) n is the set of \(x \in r\) such that \(\operatorname{ad}_{\mathfrak{g}} x\) is nilpotent; (4) n is the set of \(x \in r\) such that \(\operatorname{ad}_{\mathfrak{r}} x\) is nilpotent; (5) n is the largest ideal of g such that, for all \(x \in n\) , \(\operatorname{ad}_{\mathfrak{g}} x\) is nilpotent; (6) n is the set of \(x \in g\) such that \(\operatorname{ad}_{\mathfrak{g}} x\) belongs to the radical of the associative algebra generated by 1 and the \(\operatorname{ad}_{\mathfrak{g}} y\,(y \in g)\) .
\end{enumerate}

\textbf{(VIII) Characterizations of s:}

\begin{enumerate}
\def\labelenumi{(\arabic{enumi})}
\tightlist
\item
  \(\mathfrak{s}\) is the intersection of the kernels of the finite-dimensional simple representations of \(\mathfrak{g}\) ; (2) \(\mathfrak{s}\) is the smallest of the kernels of the finite-dimensional semi-simple representations of \(\mathfrak{g}\) ; (3) \(\mathfrak{s}\) is the intersection of the largest nilpotency ideals of the finite-dimensional representations of \(\mathfrak{g}\) ; (4) \(\mathfrak{s}\) is the smallest ideal of \(\mathfrak{g}\) such that \(\mathfrak{g} / \mathfrak{s}\) is reductive; (5) \(\mathfrak{s} = \mathfrak{r} \cap \mathcal{D}\mathfrak{g}\) ; (6) \(\mathfrak{s} = [\mathfrak{r}, \mathfrak{g}]\) ; (7) \(\mathfrak{s}\) is the intersection of the orthogonals of \(\mathfrak{g}\) relative to the bilinear forms associated with the finite-dimensional representations of \(\mathfrak{g}\) .
\end{enumerate}


\specialchapter{Nicolas Bourbaki, Elements of Mathematics}

English translations published:

General Topology, Part I

General Topology, Part II

Theory of Sets

Commutative Algebra

Algebra, Part I

Lie Groups and Lie Algebras, Part I

\textbf{General Topology, Part I}

\textbf{CHAPTER I. TOPOLOGICAL STRUCTURES}

§ 1. Open sets, neighbourhoods, closed sets

§ 2. Continuous functions

§ 3. Subspaces, quotient spaces

§ 4. Product of topological spaces

§ 5. Open mappings and closed mappings

§ 6. Filters

§ 7. Limits

§ 8. Hausdorff spaces and regular spaces

§ 9. Compact spaces and locally compact spaces

§ 10. Proper mappings

§ 11. Connectedness

Exercises

\textbf{CHAPTER II. UNIFORM STRUCTURES}

§ 1. Uniform spaces

§ 2. Uniformly continuous functions

§ 3. Complete spaces

§ 4. Relations between uniform spaces and compact spaces

\textbf{CHAPTER III. TOPOLOGICAL GROUPS}

§ 1. Topologies on groups

§ 2. Subgroups, quotient groups, homomorphisms, homogeneous spaces, product groups

§ 3. Uniform structures on groups

§ 4. Groups operating properly on a topological space; compactness in topological groups and spaces with operators

§ 5. Infinite sums in commutative groups

§ 6. Topological groups with operators; topological rings, division rings and fields

\textbf{§ 7. Inverse limits of topological groups and rings}

\textbf{CHAPTER IV. REAL NUMBERS}

§ 1. Definition of real numbers

§ 2. Fundamental topological properties of the real line

§ 3. The field of real numbers

§ 4. The extended real line

§ 5. Real-valued functions

§ 6. Continuous and semi-continuous real-valued functions

§ 7. Infinite sums and products of real numbers

§ 8. Usual expansions of real numbers; the power of R

Exercises

CHAPTER V. ONE-PARAMETER GROUPS § 1. Subgroups and quotient groups of R § 2. Measurement of magnitudes § 3. Topological characterization of the groups R and T § 4. Exponentials and logarithms Exercises

CHAPTER VI. REAL NUMBER SPACES AND PROJECTIVE SPACES § 1. Real number space R \(^{n}\) § 2. Euclidean distance, balls and spheres § 3. Real projective spaces Exercises

CHAPTER VII. THE ADDITIVE GROUPS R \(^{n}\) § 1. Subgroups and quotient groups of R \(^{n}\) § 2. Continuous homomorphisms of R \(^{n}\) and its quotient groups § 3. Infinite sums in the groups R \(^{n}\) Exercises

CHAPTER VIII. COMPLEX NUMBERS § 1. Complex numbers, quaternions § 2. Angular measure, trigonometric functions § 3. Infinite sums and products of complex numbers § 4. Complex number spaces and projective spaces Exercises

CHAPTER IX. USE OF REAL NUMBERS IN GENERAL TOPOLOGY § 1. Generation of a uniformity by a family of pseudometrics; uniformizable spaces § 2. Metric spaces and metrizable spaces § 3. Metrizable groups, valued fields, normed spaces and algebras § 4. Normal spaces § 5. Baire spaces § 6. Polish spaces, Souslin spaces, Borel sets Appendix: Infinite products in normed algebras Exercises

CHAPTER X. FUNCTION SPACES § 1. The uniformity of S-convergence § 2. Equicontinuous sets § 3. Special function spaces § 4. Approximation of continuous real-valued functions Exercises

\textbf{General Topology, Part II}

\textbf{Theory of Sets}

\textbf{CHAPTER I. DESCRIPTION OF FORMAL MATHEMATICS}

§ 1. Terms and relations

§ 2. Theorems

§ 3. Logical theories

§ 4. Quantified theories

§ 5. Equalitarian theories

Appendix. Characterization of terms and relations Exercises

CHAPTER II. THEORY OF SETS

§ 1. Collectivizing relations

§ 2. Ordered pairs

§ 3. Correspondences

§ 4. Union and intersection of a family of sets

§ 5. Product of a family of sets

§ 6. Equivalence relations

Exercises

CHAPTER III. ORDERED SETS, CARDINALS, INTEGERS

§ 1. Order relations. Ordered sets

§ 2. Well-ordered sets

§ 3. Equipotent sets. Cardinals

§ 4. Natural integers. Finite sets

§ 5. Properties of integers

§ 6. Infinite sets

§ 7. Inverse limits and direct limits

Exercises

CHAPTER IV. STRUCTURES

§ 1. Structures and isomorphisms

§ 2. Morphisms and derived structures

§ 3. Universal mappings

Exercises

SUMMARY OF RESULTS

Introduction

§ 1. Elements and subsets of a set

§ 2. Functions

§ 3. Products of sets

§ 4. Union, intersection, product of a family of sets

§ 5. Equivalence relations and quotient sets

§ 6. Ordered sets

§ 7. Powers. Countable sets

§ 8. Scales of sets. Structures

\textbf{Commutative Algebra}

CHAPTER I. FLAT MODULES § 1. Diagrams and exact sequences § 2. Flat modules § 3. Faithfully flat modules § 4. Flat modules and ``Tor'' functors Exercises

CHAPTER II. LOCALIZATION § 1. Prime ideals § 2. Rings and modules of fractions § 3. Local rings. Passage from the local to the global § 4. Spectra of rings and supports of modules § 5. Finitely generated projective modules. Invertible fractional ideals Exercises

CHAPTER III. GRADUATIONS, FILTRATIONS AND TOPOLOGIES § 1. Finitely generated graded algebras § 2. General results on filtered rings and modules § 3. m-adic topologies on Noetherian rings § 4. Lifting in complete rings § 5. Flatness properties of filtered modules Exercises

CHAPTER IV. ASSOCIATED PRIME IDEALS AND PRIMARY DECOMPOSITION § 1. Prime ideals associated with a module § 2. Primary decomposition § 3. Primary decomposition in graded modules Exercises

CHAPTER V. INTEGERS § 1. Notion of an integral element § 2. The lift of prime ideals § 3. Finitely generated algebras over a field Exercises

CHAPTER VI. VALUATIONS § 1. Valuation rings § 2. Places § 3. Valuations § 4. The height of a valuation § 5. The topology defined by a valuation § 6. Absolute values § 7. Approximation theorem § 8. Extensions of a valuation to an algebraic extension § 9. Application: locally compact fields § 10. Extensions of a valuation to a transcendental extension Exercises

\textbf{COMMUTATIVE ALGEBRA}

CHAPTER VII. DIVISORS

§ 1. Krull domains

§ 2. Dedekind domains

§ 3. Factorial domains

§ 4. Modules over integrally closed Noetherian domains Exercises

\textbf{Algebra, Part I}

\textbf{CHAPTER I. ALGEBRAIC STRUCTURES}

§ 1. Laws of composition; associativity; commutativity

§ 2. Identity element; cancellable elements; invertible elements

§ 3. Actions

§ 4. Groups and groups with operators

§ 5. Groups operating on a set

§ 6. Extensions, solvable groups, nilpotent groups

§ 7. Free monoids, free groups

§ 8. Rings

§ 9. Fields

§ 10. Inverse and direct limits

Exercises

\textbf{CHAPTER II. LINEAR ALGEBRA}

§ 1. Modules

§ 2. Modules of linear mappings. Duality

§ 3. Tensor products

§ 4. Relations between tensor products and homomorphism modules

§ 5. Extension of the ring of scalars

§ 6. Inverse and direct limits of modules

§ 7. Vector spaces

§ 8. Restruction of the field of scalars in vector spaces

§ 9. Affine spaces and projective spaces

§ 10. Matrices

§ 11. Graded modules and rings

Appendix. Pseudomodules

Exercises

\textbf{CHAPTER III. TENSOR ALGEBRAS, EXTERIOR ALGEBRAS, SYMMETRIC ALGEBRAS}

§ 1. Algebras

§ 2. Examples of algebras

§ 3. Graded algebras

§ 4. Tensor products of algebras

§ 5. Tensor algebra. Tensors

§ 6. Symmetric algebras

§ 7. Exterior algebras

§ 8. Determinants

§ 9. Norms and traces

§ 10. Derivations

§ 11. Cogebras, products of multilinear forms, inner products and duality Appendix. Alternative algebras. Octonions

Exercises
